%\aaunnumberedsection{ВВЕДЕНИЕ}{sec:intro}
%
%Параллельные вычисления позволяют повысить производительность обработки данных и широко используются. Алгоритм Z-буфера прост в реализации и часто применим на практике~\cite{rogers}.
%
%Целью данной работы является исследование параллельной и непараллельной реализаций алгоритма Z-буфера.
%
%\label{sec:goals}
%Для достижения поставленной цели необходимо выполнить следующие задачи:
%\begin{itemize}[label=---]
%	\item написать программу, реализующую параллельный и непараллельный алгоритм Z-буфера;
%	\item проверить работоспособность программы, выполнив тестирование;
%	\item замерить процессорное время работы алгоритмов;
%	\item провести анализ полученных результатов.
%\end{itemize}
%
%\aasection{Входные и выходные данные}{sec:input-output}
%
%Входным данным для программы является файл с описанием многоугольников. Этот файл должен содержать количество N многоугольников (целое число). Далее в файле должны быть перечислены N многоугольников в формате: кол-во вершин M в текущем многоугольнике (целое число), а также M самих вершин этого многоугольника в формате x y z.
%
%\aasection{Преобразование входных данных в выходные}{sec:algorithm}
%Программа считывает многоугольники из файла и выводит их на экран в порядке, видимом пользователю (применяется алгоритм Z-буфера).
%
%\aasection{Примеры работы программы}{sec:demo}
%
%На рисунке~\ref{img:workExample} представлен пример работы программы.
%
%\FloatBarrier
%\includeimage
%{workExample} % Имя файла без расширения (файл должен быть расположен в директории inc/img/)
%{f} % Обтекание (без обтекания)
%{h} % Положение рисунка (см. figure из пакета float)
%{\textwidth} % Ширина рисунка
%{Пример работы программы} % Подпись рисунка
%\FloatBarrier
%
%
%\aasection{Тестирование}{sec:tests}
%
%В таблице~\ref{tbl:tests} представлены функциональные тесты для разработанного ПО. Все тесты пройдены успешно. Размер буфера кадра для тестирования~---~4x4. 
%
%Пусть:
%\begin{itemize}[label=---]
%	\item w~---~белый цвет;
%	\item r~---~красный цвет;
%	\item b~---~синий цвет.
%\end{itemize}
%А перед перечислением вершин многоугольника также указан его цвет. Тогда:
%
%\begin{longtable}{|p{.14\textwidth - 2\tabcolsep}|p{.36\textwidth - 2\tabcolsep}|p{.24\textwidth - 2\tabcolsep}|p{.23\textwidth - 2\tabcolsep}|}
%    \caption{Функциональные тесты}\label{tbl:tests} \\\hline
%    № теста & Входные данные & Полученные выходные данные & Ожидаемые выходные данные                                          \\\hline
%    \endfirsthead
%    \caption{Функциональные тесты (продолжение)} \\\hline
%    № теста & Входные данные & Полученные выходные данные  & Ожидаемые выходные данные                                                 \\\hline
%    \endhead
%    \endfoot
%    1                                           & 1 \newline 3 r\newline [0 0 1] [0 4 1] [4 4 1] & [[r, w, w, w], \newline  [r, r, w, w],\newline [r, r, r, w],\newline [r, r, r, r]] & [[r, w, w, w], \newline [r, r, w, w], \newline[r, r, r, w],\newline [r, r, r, r]] \\\hline
%    2                                           & 2 \newline 3 r \newline [0 0 1] [0 4 1] [4 4 1] \newline 3 b \newline [0 0 0] [0 4 0] [4 4 0] & [[r, w, w, w], \newline  [r, r, w, w],\newline [r, r, r, w],\newline [r, r, r, r]] & [[r, w, w, w], \newline [r, r, w, w], \newline[r, r, r, w],\newline [r, r, r, r]] \\\hline
%    3                                           & 2 \newline 4 r \newline [0 0 1] [0 4 2] [4 4 1] [4 0 0] \newline 4 b \newline [0 0 1] [0 4 0] [4 4 1] [4 0 2] & [[r, b, b, b], \newline  [r, r, b, b],\newline [r, r, r, b],\newline [r, r, r, r]] & [[r, b, b, b], \newline  [r, r, b, b],\newline [r, r, r, b],\newline [r, r, r, r]]  \\\hline
%    \end{longtable}
%
%\aasection{Описание исследования}{sec:study}
%В ходе исследования требуется замерить время работы алгоритма.
%\newpage
%
%В таблице~\ref{tbl:time} приведены полученные результаты замеров.
%
%Технические характеристики устройства, на котором выполнялись замеры времени, представлены далее.
%
%\begin{enumerate}
%	\item Процессор	AMD Ryzen 5 5500u ЦПУ 2592 МГц, ядер: 8, логических процессоров: 12.
%	\item Оперативная память: 16 ГБ.
%	\item Операционная система: Windows 10.
%\end{enumerate}
%
%При замерах времени ноутбук был включен в сеть электропитания. Запущены были только системные приложения.
%
%Замеры проводились при одинаковых входных данных с количеством многоугольников~---~5.
%\begin{longtable}{|c|c|}
%	\caption{Результаты замеров времени}\label{tbl:time} \\\hline
%	Алгоритм & Полученное время, мс за 1 кадр  \\\hline
%	\endfirsthead
%	Z-буфер, нет потоков & 11 \\\hline
%	Z-буфер, 1 поток & 13 \\\hline
%	Z-буфер, 2 потока & 16 \\\hline
%	Z-буфер, 4 потока & 20 \\\hline
%	Z-буфер, 8 потоков & 36 \\\hline
%	Z-буфер, 16 потоков & 42 \\\hline
%	Z-буфер, 32 потока & 53 \\\hline
%	Z-буфер, 64 потока & 71 \\\hline
%\end{longtable}
%
%В ходе исследования получено, что реализация параллельного алгоритма Z-буфера работает медленнее, чем непараллельная.
%
%\aaunnumberedsection{ЗАКЛЮЧЕНИЕ}{sec:outro}
%
%Поставленная цель была достигнута.
%Для достижения поставленной цели были выполнены все поставленные задачи:
%\begin{itemize}[label=---]
%	\item написана программа, реализующая параллельный и непараллельный алгоритм Z-буфера;
%	\item проверена работоспособность программы, выполнено тестирование;
%	\item произведены замеры процессорного времени работы алгоритмов;
%	\item проведён анализ полученных результатов.
%\end{itemize}
